\section{Analisi dei dati e classificatori utilizzati}

In questa sezione verrà riportato lo studio effettuato sul dataset.
% Ogni sottosezione tratterà di uno specifico classificatore, dei suoi risultati e degli esperimenti fatti per migliorarne il modello.

\subsection{Considerazioni sugli attributi}

Studiando la composizione dei dati per selezionare le principali tecniche da applicare,
la prima cosa che risulta visibile all'apertura del dataset è sicuramente il suo sbilanciamento nelle classi di output:
come mostrato in~\Cref{fig:classes}, su 100 pazienti in esame, solo 12 risultano compromessi.

\begin{figure}[H]
  \centering
  \includegraphics[width=0.85\textwidth]{fig/classes.eps}%
  \caption{%
    L'istogramma rappresenta i valori assunti dall'attributo classe:
    la colonna blu rappresenta le diagnosi di fertilità normale,
    quella rossa gli esami che evidenziano potenziali problematiche.
  }%
  \label{fig:classes}
\end{figure}

Questo disequilibrio tra le classi mette in difficoltà i modelli più semplici;
al fine di ottenere risultati validi, si è messo alla prova diverse configurazioni parametriche.

\todo[inline]{Aggiungere maggiori dettagli per ciascun attributo}

\subsection{Misure di performance dei classificatori}

\todo[inline]{Migliorare quanto segue}

Prima di partire con i vari dati, occorre però stabilire i criteri di misura delle performance dei vari classificatori:
a tale scopo ho tenuto conto in primo luogo della Confusion Matrix prodotta da ogni modello, successivamente dei valori di Precision e Recall per le classi in gioco, attribuendo un maggiore peso ai valori relativi alla classe ``O'' rispetto a ``N'' e infine del valore della ROC\@.
Come modalità di test ho scelto di applicare in ogni situazione la K-Fold Cross-Validation piuttosto che lo splitting statico dei dati in Training e Test Set, la motivazione dietro a questa scelta è l'esiguo numero di dati presenti nel dataset a cui si presta meglio una modalità come il K-Fold Cross-Validation,
in quanto pensata appositamente per dataset di piccole dimensioni.

\todo[inline]{Manca qualcosa?}

\subsection{Decision Tree}

La prima tecnica di classificazione testata è stata il Decision Tree.
In particolare, si è scelto di usare \texttt{J48}, implementazione dell'algoritmo \emph{C4.5};
esso realizza un albero binario a cui possono essere applicate tecniche di \emph{post pruning}
per ridurre le dimensioni dell'albero e l'errore commesso dal classificatore.

In primo luogo si è provato a generare un modello applicando i parametri standard impostati da Weka, ottenendo così i risultati riportati in~\Cref{fig:j48}.

\begin{figure}[H]
  \centering
  \adjincludegraphics[width=\linewidth,trim=0 0 0 {.3\height},clip]{fig/j48.eps}%
  \caption{Output di \texttt{J48} eseguito tramite Weka}%
  \label{fig:j48}
\end{figure}

Com'è possibile notare, il classificatore costruisce un modello poco valido:
per quanto \emph{accuracy}, \emph{precision} e \emph{recall} per la classe ``\texttt{N}'' abbiano valori piuttosto elevati,
i corrispondenti per ``\texttt{O}'' risultano tendenti a \(0\).
Anche il valore di \emph{ROC}, indicatore principale della bontà di un modello di classificazione, risulta estremamente basso.

Osservando la struttura dell'albero costruito, è piuttosto chiaro il motivo di valori tanto bassi:
l'albero infatti è costruito da un'unica foglia i cui valori sono assegnati alla classe ``\texttt{N}'', causando così la situazione sopra descritta.

Ovviamente questo modello non risultava di una qualità soddisfacente, probabilmente a causa del post pruning,
il quale, in presenza di un tale sbilanciamento nella frequenza delle classi, tende a ridurre l'albero di classificazione a singola foglia.
Per porre rimedio a questa cosa si è provato ad aumentare e diminuire \improvement{magari mettere il valore usato? 0.25} il parametro \texttt{confidenceFactor}, responsabile di accentuare o smussare le operazioni di pruning;
anche però variando questo parametro i risultati ottenuti non sono stati particolarmente differenti.

La situazione è decisamente migliorata disattivando il pruning (tramite parametro \texttt{unpruned}):
questa scelta ha finalmente prodotto un risultato notevolmente differente dai precedenti, come è possibile vedere in~\Cref{fig:j48-unpruned}.

\begin{figure}[H]
  \centering
  \begin{subfigure}{0.35\textwidth}
    \centering
    \adjincludegraphics[width=\linewidth,trim=0 {.37\height} {.55\width} {.25\height},clip]{fig/j48-unpruned.eps}%
    \label{fig:j48-unpruned:tree}
  \end{subfigure}
  \hfill
  \begin{subfigure}{0.6\textwidth}
    \centering
    \adjincludegraphics[width=\linewidth,trim=0 0 {.16\width} {.65\height},clip]{fig/j48-unpruned.eps}%
    \label{subfig:j48-unpruned:result}
  \end{subfigure}
  \caption{Albero di classificazione e output ottenuto con \texttt{J48} \emph{unpruned}}%
  \label{fig:j48-unpruned}
\end{figure}

Il risultato del classificatore è sicuramente migliore rispetto alla sua versione \emph{pruned}, tuttavia rimane ancora abbastanza insoddisfacente.

Sperimentando ulteriormente con la versione \emph{unpruned} di \texttt{J48},
abbassando il numero minimo di elementi per foglia a \(1\) il valore di \emph{recall} per la classe \texttt{O} e di \emph{True Positive} si alzava,
ma il valore di \emph{ROC} e \emph{precision} si abbassavano.

Si può dunque concludere che il classificatore \texttt{J48} non garantisce risultati ottimali sopra questo dataset, soprattutto a causa del grande sbilanciamento tra le distribuzioni delle classi.

Disabilitando il post pruning e abbassando il numero minimo di elementi per foglia a \(1\), si può classificare correttamente poco meno della metà degli elementi appartenente a \texttt{O},
classe più difficile da individuare ma allo stesso tempo più rilevante rispetto al dominio dei dati del sistema.

\subsection{Classificatori Bayesiani: \emph{Naïve Bayes}}

Per tentare un approccio differente, ho provato a costruire un modello basato sui classificatori bayesiani,
approccio che ignora le relazioni tra gli attributi in favore della frequenza di certi valori rispetto alle classi in gioco.

Il classificatore specifico scelto per il dataset è stato \emph{Naïve Bayes} nella sua implementazione fornita da Weka (chiamata appunto \texttt{NaiveBayes});
costruendo il modello utilizzando i parametri standard impostati dalla workbench, sono stati ottenuti i risultati mostrati in~\Cref{fig:bayes}.

\begin{figure}[H]
  \centering
  \adjincludegraphics[width=\linewidth,trim=0 0 0 {.76\height},clip]{fig/bayes.eps}%
  \caption{Output di \texttt{NaiveBayes} eseguito tramite Weka}%
  \label{fig:bayes}
\end{figure}

Analizzando i risultati di questo classificatore, si possono fare diverse considerazioni:
partendo dall'osservazione della \emph{confusion matrix} è possibile vedere come praticamente tutti i dati vengano assegnati alla classe ``\texttt{N}'', mentre ad ``\texttt{O}'' viene assegnato un solo elemento;
ciò è ulteriormente confermato dai valori i \emph{precision} e \emph{recall} per la classe ``\texttt{O}'', che risultano essere \(0\).

Infine, osservando \emph{ROC} si ha la conferma definitiva che il modello costruito sia decisamente inefficiente e poco performante.

Le ragioni dietro alla poca efficienza del modello bayesiano per questo problema possono essere ricondotte alla mancanza dei valori di distribuzione delle probabilità reali per le classi e alla potenziale mancanza di indipendenza statistica tra attributi.
Questi fattori, combinati allo sbilanciamento dei dati nel dataset, da cui vengono estratte le distribuzioni di probabilità necessarie al classificatore in mancanza di quelle reali, rendono inefficace il modello bayesiano su questo dataset.

\subsection{Classificatori Lazy: \emph{k-NN}}

Altra famiglia utilizzata è stata quella dei \emph{classificatori Instance-based} o \emph{Lazy}.
La particolarità di queste tecniche è la totale assenza di modello, sostituito da procedure di natura induttiva.

Tra tutti i classificatori disponibili in questa famiglia, si è scelto di utilizzare \texttt{IBk}, implementazione in Weka dell'algoritmo \emph{k-NN};
per questo classificatore è possibile impostare diversi parametri, come il numero di vicini da considerare e la metrica da utilizzare nel calcolo della distanza.

\subsubsection{Utilizzando gli attributi nominali come tali}

Applicando i parametri default di Weka, ovvero \texttt{KNN=1} e \emph{distanza euclidea}, si ottengono i risultati riportati in~\Cref{fig:ibk:1}.

\begin{figure}[H]
  \centering
  \adjincludegraphics[width=\linewidth,trim=0 0 {.285\width} {.32\height},clip]{fig/ibk.eps}%
  \caption{Output di \texttt{IBk} eseguito tramite Weka}%
  \label{fig:ibk:1}
\end{figure}

I risultati ottenuti da questo modello risultano non particolarmente buoni, tuttavia i valori di \emph{Confusion Matrix}, \emph{ROC}, \emph{Precision} e \emph{Recall} per la classe ``\texttt{O}'' sono comunque migliori rispetto a quelli ottenuti nei tentativi precedenti.
Si è ritenuto dunque che il modello possa essere migliorato utilizzando valori differenti per i parametri sopra descritti.

Mantenendo la \emph{distanza euclidea} come metrica, si è provato a variare il valore di K tra i numeri dispari compresi tra \(3\) e \(11\);
l'unico risultato rilevante è con \texttt{KNN=9}, per il quale si raggiunge la massimizzazione del valore di \emph{ROC} (\Cref{fig:ibk:9}).

%     % \begin{figure}[H]
%     %   \centering
%     %   \adjincludegraphics[width=\linewidth,trim=0 0 {.285\width} {.32\height},clip]{fig/ibk5.eps}%
%     %   \caption{Output di \texttt{IBk} con \texttt{KNN=5}}%
%     %   \label{fig:ibk:5}
%     % \end{figure}
%   \item
%     con \texttt{KNN=3} si ottengono i massimi valori di \emph{Precision} e \emph{Recall} per la classe ``\texttt{O}'' (\Cref{fig:ibk:3}).
%     % \begin{figure}[H]
%     %   \centering
%     %   \adjincludegraphics[width=\linewidth,trim=0 0 {.285\width} {.32\height},clip]{fig/ibk3.eps}%
%     %   \caption{Output di \texttt{IBk} con \texttt{KNN=3}}%
%     %   \label{fig:ibk:3}
%     % \end{figure}
% \end{itemize}

% \begin{figure}[H]
%   \centering
%   \adjincludegraphics[width=\linewidth,trim=0 0 {.285\width} {.32\height},clip]{fig/ibk3.eps}%
%   \caption{Output di \texttt{IBk} con \texttt{KNN=3}}%
%   \label{fig:ibk:3}
% \end{figure}

\begin{figure}[H]
  \centering
  \adjincludegraphics[width=\linewidth,trim=0 0 {.285\width} {.32\height},clip]{fig/ibk9.eps}%
  \caption{Output di \texttt{IBk} con \texttt{KNN=9}}%
  \label{fig:ibk:9}
\end{figure}

Di fatto, anche cambiando il numero dei vicini da considerare si può vedere che il modello arranca ancora nel riconoscimento degli elementi della classe ``\texttt{O}''.

\subsubsection{Utilizzando gli attributi come numerici}

Applicando gli stessi parametri per configurare il numero di vicini (tutti i numeri dispari nel range \(1\)-\(11\), con \emph{distanza euclidea}), i risultati più rilevanti sono con \texttt{KNN=3} e \texttt{KNN=5}:
infatti, se ad esempio con \texttt{KNN=1} (opzione di default) i risultati sono leggermente peggiori rispetto all'utilizzo dei dati in formato categorico, i seguenti risultati sono invece maggiormente rilevanti:

\begin{itemize}
  \item con \texttt{KNN=3} (\Cref{fig:ibk:3-num}), si ottengono valori di \emph{precision} e di \emph{recall} per la classe ``\texttt{O}'' migliori;
  \item con \texttt{KNN=5} (\Cref{fig:ibk:5-num}), si raggiunge la massimizzazione di \emph{ROC}.
\end{itemize}

% \begin{figure}[htbp]
%   \centering
%   \adjincludegraphics[width=\linewidth,trim=0 {.07\width} {.285\width} {.32\height},clip]{fig/ibk.eps}%
%   \caption{Output di \texttt{IBk} con \texttt{KNN=1} utilizzando gli attributi come numerici}%
%   \label{fig:ibk:1-num}
% \end{figure}

\begin{figure}[H]
  \centering
  \adjincludegraphics[width=\linewidth,trim=0 {.07\width} {.285\width} {.32\height},clip]{fig/ibk3-num.eps}%
  \caption{Output di \texttt{IBk} con \texttt{KNN=3} utilizzando gli attributi come numerici}%
  \label{fig:ibk:3-num}
\end{figure}

\begin{figure}[H]
  \centering
  \adjincludegraphics[width=\linewidth,trim=0 {.07\width} {.285\width} {.32\height},clip]{fig/ibk5-num.eps}%
  \caption{Output di \texttt{IBk} con \texttt{KNN=5} utilizzando gli attributi come numerici}%
  \label{fig:ibk:5-num}
\end{figure}

Al netto di ciò, si può dire che l'utilizzo degli attributi come numerici permette di avere effettivamente risultati migliori,
tuttavia il modello arranca ancora nel riconoscimento degli elementi della classe ``\texttt{O}''.
Le ragioni di questo comportamento sono molto probabilmente non solo dovute allo sbilanciamento nelle classi di output,
ma anche alla mancanza di una separazione netta tra i dati a livello di valori degli attributi.

Ritenendo di poter essere sulla buona strada, si è provato a effettuare un cambio di metrica per misurare la distanza tra punti (ovvero i dati) nello spazio:
tra le metriche a disposizione nella workbench, si è provata la \emph{distanza di Čebyšëv} (\texttt{ChebyshevDistance} in Weka), che misura la distanza tra due vettori come la maggiore tra le differenze dei singoli attributi.
Nelle~\Cref{fig:ibk:3-num-che,fig:ibk:5-num-che} sono riportati i risultati per \texttt{KNN=3} e \texttt{KNN=5}, che però non mostrano miglioramenti tali da poter considerare l'algoritmo particolarmente utile per il caso di studio.

\begin{figure}[H]
  \centering
  \adjincludegraphics[width=0.9\linewidth,trim=0 {.1\width} {.285\width} {.32\height},clip]{fig/ibk3-num-che.eps}%
  \caption{Output di \texttt{IBk} con \texttt{KNN=3} utilizzando attributi numerici e \emph{distanza di Čebyšëv}}%
  \label{fig:ibk:3-num-che}
\end{figure}

\begin{figure}[H]
  \centering
  \adjincludegraphics[width=0.9\linewidth,trim=0 {.1\width} {.285\width} {.32\height},clip]{fig/ibk5-num-che.eps}%
  \caption{Output di \texttt{IBk} con \texttt{KNN=5} utilizzando attributi numerici e \emph{distanza di Čebyšëv}}%
  \label{fig:ibk:5-num-che}
\end{figure}

\subsection{Multi-Classificatori}

\subsubsection{Approccio Cost-Sensitive}

\subsubsection{Approccio Boosting}

\subsection{Rules}
